\section{Formule Normale de Frobenius}
\subsection{Déterminer la forme normale de Frobenius}
On a que pour une matrice $A \in \C^{n \times n}$ alors il existe une matrice $P \in \C^{n \times n}$ telle que $A = P J P^{-1}$ avec $J \in \C^{n \times n}$ une forme normale de Jordan. Or si on se place dans un coprs $\K$ quelconque, ce résultat est vrai seulement si le polynôme caractéristique $\chi_A$ est scindé, ce qui n'est pas toujours le cas, comme par exemple si $A \in \R$. 

\newdef{La \emph{matrice compagnon} d'un polynôme $f(x) = a_0 + a_1 x + \ldots a_{d-1}x^{d-1} + x^d \in \K[X]$ est la matrice
	\[
	C_f = \begin{pmatrix}   0 &   &  &  & - a_0 \\
							1 & 0 &  &  & -a_1 \\
							  & \ddots & \ddots &  & \vdots \\
							  &        &  1 & 0  & -a_{d-2} \\
							  &			&       & 1 & -a_{d-1}
		\end{pmatrix}
		\in \K^{d \times d}
	\]
	Il est alors facile de voir que $\chi_{C_f} = f$.
}{}

\newthm{Soit $A \in \K^{n \times n}$, alors il existe une matrice inversible $T \in \K^{n \times n}$ telle que
	\[
	A = T \begin{pmatrix} C_{f_1} &  & \\
							      & \ddots &  \\
							      &        & C_{f_k}
	\end{pmatrix} T^{-1}
	\]
	où chaque $C_{f_i}$ est une matrice compagnon avec $f_{i+1} \mid f_i$ pour tout $i = 1, \ldots , k-1$.
}{}

\prf{
On commence par calculer le polynôme minimal $\mu_A (x) = a_0 + a_1 x + \ldots + a_{d-1} x^{d-1} + x^d \in \K[X]$ ainsi qu'un vecteur $v \in \K^n$ tel que $p_v(x) = \mu_A(x)$. Ainsi les vecteurs
\[
v, Av, \ldots, A^{d-1} v 
\]
sont linéairement indépendant, car sinon il existerait une combinaison linéaire non trivial nulle, ce qui contredirait la minimalité de $p_v(x)$. De plus
\[
A^d v = -a_0 v - a_1 Av - \ldots - a_{d-1} A^{d-1} v
\]
soit $J = \{j_1, \ldots, j_{n-d}\} \subset \entiers$ un ensemble d'indice tell que
\[
\mathcal B = \{v, Av, \ldots, A^{d-1}v, e_{j_1}, \ldots, e_{j_{n-d}}\}
\]
est une base de $\K^n$ et soit $T \in \K^{n \times n}$ la matrice dont les colonnes sont les éléments de $\mathcal B$ dans cet ordre. Par construction, la matrice $A$ est de la forme
\[
A_{\mathcal B} = \begin{pmatrix} C_{\mu_A} & A_1 \\ 0  & A_2 \end{pmatrix} \in \K^{n \times n}
\]
où $\binom{A_1}{A_2} \in \K^{n \times ({n-d})}$. Par la formule de changement de base, on a
\[
A_{\mathcal B} = T^{-1} A T
\]
Maintenant le but est d'utiliser successivement des opérations élémentaires sur les lignes et colonnes pour que la matrice dans le coin supérieure droit (initialement $A_1$) devienne nulle.  Plus précisément, $A_{\mathcal B}$ ressemble à
\[
A_{\mathcal B}=
\begin{pmatrix}
	0      &        &        &        & -a_0     & b_{1,1}   & \cdots & b_{1,n-d} \\
	1      & 0      &        &        & \vdots   & b_{2,1}   & \cdots & b_{2,n-d} \\
	& \ddots & \ddots &        & \vdots   & \vdots    &        & \vdots \\
	&        & 1      & 0      & -a_{d-2} & \vdots    & \cdots & \vdots \\
	&        &        & 1      & -a_{d-1} & b_{d,1}   & \cdots & b_{d,n-d} \\
	      &  &  &       &         & b_{d+1,1} & \cdots & b_{d+1,n-d} \\
 &        &        &  &    & \vdots    &        & \vdots \\
	      &  &  &       &         & b_{n,1}   & \cdots & b_{n,n-d}
\end{pmatrix}
\in K^{n\times n}.
\]
On peut donc faire disparaître la $d$-ème ligne, donc on applique les opération élementaire sur les colonnes  
\[
C_d \leftarrow C_d - b_{d,1} C_{d-1}, \quad \ldots \quad C_n \leftarrow C_n - b_{d, n-d} C_{d-1}
\]
ainsi que ses inverses correspondants sur les lignes
\[
L_{d-1} \leftarrow L_{d-1} + b_{d,1} L_d, \quad \ldots \quad L_{d-1} \leftarrow L_{d-1} + b_{d, n-d}
\]
ainsi on recommence cela à chaque ligne, pour finalement arriver à 
\[
B' =
\begin{pmatrix}
	0      &        &        &        & -a_0     &    &  &  \\
	1      & 0      &        &        & \vdots   &    &  &  \\
	& \ddots & \ddots &        & \vdots   &     &        &  \\
	&        & 1      & 0      & -a_{d-2} &     &  &  \\
	&        &        & 1      & -a_{d-1} &    &  &  \\
	&  &  &       &         & b_{d+1,1} & \cdots & b_{d+1,n-d} \\
	&        &        &  &    & \vdots    &        & \vdots \\
	&  &  &       &         & b_{n,1}   & \cdots & b_{n,n-d}
\end{pmatrix}
\in K^{n\times n}.
\]
Ainsi après cette procédure, on a transformer la matrice dans la forme
\[
\begin{pmatrix} C_{\mu_A} &  \\  & A' \end{pmatrix} \in \K^{n \times n}
\]
le polynôme minimale de $A'$ divise bien $\mu_A$. Par récurrence, on a bien montrer le résultat voulu.
}

\subsection{Utilisation de la forme normale de Frobenius}
Si on considère une matrice compagnon $C_f \in \K^{d \times d}$ alors on sait que
\[
f(C_f) = 0
\]
Donc toute puissance élevée de $C_f$ peut être écrite comme combinaison linéaire de 
\[
I, C_f, \ldots, C_f^{d-1}
\]
Donc si j'ai un polynôme $p \in \K[X]$, on a
\[
p(x) = q(x) f(x) + r(x)
\]
avec $\deg r < \deg f$ et donc par évaluation en $C_f$ on a
\[
p(C_f) = r(C_f)
\]
Ainsi pour une matrice $A \in \K$ et $p \in \K[X]$ on a
\[
p(A) \sim \begin{pmatrix} p(C_{f_1}) &  & \\
							& \ddots & \\
							&  & p(C_{f_k})
		\end{pmatrix}
		= 
		\begin{pmatrix} r_1(C_{f_1}) &  & \\
			& \ddots & \\
			&  & r_k(C_{f_k})
		\end{pmatrix}
\]
Ainsi il est plus simple de calculer $p(A)$ car $r_i$ est de degré petit. Et donc calculer des grandes puissances de matrices revient à calculer des petites puissance de matrice compagnon, ce qui est bien plus simple.

Et donc dans le cas où on souhaite calculer $e^A$, on a que l'exponentielle devient un simple polynôme en $A$
\[
e^A = \alpha_0 I + \alpha_1 A + \ldots + \alpha_{d-1} A^{d-1}
\]
\newthm{Deux matrices sont semblables si et seulement si elles sont les mêmes forme normale de Frobenius}{}
